% !TEX root = ../../main.tex
% C3.7 - RTSP and camera simulation (translated). RFC 2326/7826/3550/6184, H.264, FFmpeg, MediaMTX.
% Matches scripts/rtsp.ps1 and workers/sources/rtsp.py (PyAV, TCP, 8 s timeouts, 3 reconnects 0.25-2 s, worker-clock timestamps).
\section{RTSP and Camera Stream Simulation}
\label{sec:3.7}

\subsection{H.264 Compressed Video}
\label{sec:3.7.1}

IP cameras compress video, since an uncompressed 1920$\times$1080 frame takes about 6~MB, or about 370~MB per second at 60 frames per second. H.264~\cite{wiegand2003h264} is the common standard and the project's format. An \emph{I-frame} only uses prediction within the frame itself, so it can be decoded independently; \emph{P-frames} and \emph{B-frames} only store motion and differences relative to reference frames, so they can only be decoded once their references are available. Decoding must therefore be done on every frame even though the system only analyses some of them (Section~\ref{sec:3.4}), and a decoder that joins a stream midway must wait for the next I-frame.

\subsection{RTSP and RTP}
\label{sec:3.7.2}

RTSP (Real Time Streaming Protocol)~\cite{rfc2326} is a \emph{control} protocol: the receiver identifies a stream by an address of the form \texttt{rtsp://<host>:<port>/<path>} and exchanges the \texttt{DESCRIBE}, \texttt{SETUP}, \texttt{PLAY} and \texttt{TEARDOWN} commands to obtain the description, set up, start and end the session. The data is carried by RTP (Real-time Transport Protocol)~\cite{rfc3550}: compressed video is split into packets with sequence numbers and timestamps; for H.264 the timestamps use a 90~kHz clock~\cite{rfc6184}. RTP packets can be sent over UDP, with low latency but possible packet loss, or interleaved on the TCP connection of the RTSP session, without loss but possibly late under congestion. RTSP 2.0~\cite{rfc7826} removed the commands used to push a stream to a server, so the project's simulation tools use version 1.0.

\subsection{Simulating Cameras with a Streaming Server}
\label{sec:3.7.3}

To test with the real protocol of IP cameras, each video is replayed as an RTSP stream acting as one camera. The MediaMTX streaming server~\cite{mediamtx2026} accepts streams pushed in and serves streams out over RTSP, each at its own path such as \texttt{/cam1}. The FFmpeg tool~\cite{ffmpeg2026docs} reads the file and pushes it to MediaMTX with three options: \texttt{-re} to send at the original frame rate, \texttt{-stream\_loop -1} to loop forever like a camera that never stops, and \texttt{-c copy} to copy the compressed data without re-encoding, so the stream carries exactly the H.264 data of the original file and costs almost no resources. Like a real camera, the stream uses the same protocol and format and does \emph{not wait} for a slow receiver; unlike one, its content repeats and it suffers almost no packet loss or disconnection.

\subsection{Receiving RTSP Streams in the System}
\label{sec:3.7.4}

The background worker reads streams with the PyAV library, which uses FFmpeg's decoding libraries, with RTP interleaved over TCP to avoid packet loss corrupting images. Connecting and waiting for data are both limited to 8 seconds; when a stream is interrupted, the system tries to reconnect up to three times with a waiting time growing from 0.25 to 2 seconds, then records the processing run as failed. Each frame is timestamped from the stream's own presentation timestamp, anchored to the machine clock at the first frame of each session, so that the time between two processed frames is the stream time between them even when the machine falls behind; a session that starts again after a reconnection or a restart of the source is anchored again, and only a frame without a usable timestamp falls back to the reception time (Section~\ref{sec:3.3.5}). Processing is much slower than the stream, so cameras share processing time as described in Section~\ref{sec:5.2.1}.
